% ========== UI EXTENSIBILITY ==========
% Custom views, dynamic layout modifications, and component registration
% Source: Symphony/Content/The Minimal IDE, The Assembling

\section*{UI Extensibility Framework}
\addcontentsline{toc}{section}{UI Extensibility Framework}

\lettrine{S}{ymphony's} UI extensibility framework represents one of the most sophisticated aspects of its generic primitives approach. Rather than providing fixed UI components or rigid layout systems, Symphony offers a flexible framework that allows extensions to create entirely custom user interfaces while maintaining consistency and performance.

\subsection*{Custom Views and Components}
\addcontentsline{toc}{subsection}{Custom Views and Components}

Extensions can create custom views that integrate seamlessly with Symphony's core interface. The UI extensibility system provides several levels of customization:

\begin{table}[h]
\centering
\begin{tabular}{@{}p{0.3\textwidth}p{0.35\textwidth}p{0.3\textwidth}@{}}
\toprule
\textbf{Customization Level} & \textbf{Capabilities} & \textbf{Use Cases} \\
\midrule
\textbf{Component Registration} & 
Custom React components, event handlers, styling & 
Status bars, toolbars, custom buttons \\
\midrule
\textbf{View Creation} & 
Full custom views, data binding, lifecycle management & 
File explorers, output panels, search results \\
\midrule
\textbf{Layout Modification} & 
Panel arrangement, window management, workspace organization & 
IDE layouts, multi-pane editors, dashboard views \\
\midrule
\textbf{Virtual DOM Bridge} & 
Direct DOM manipulation, custom rendering, performance optimization & 
Complex visualizations, real-time displays, games \\
\bottomrule
\end{tabular}
\caption{UI Extensibility Levels and Capabilities}
\end{table}

\subsubsection*{Component Registration System}

Extensions register UI components through a declarative manifest system:

\begin{lstlisting}[language=toml, caption=UI Component Registration]
[ui_components]
# Register a custom sidebar panel
[[ui_components.panels]]
id = "git-status-panel"
title = "Git Status"
location = "sidebar-left"
icon = "git-branch"
default_visible = true

# Register a custom toolbar button
[[ui_components.toolbar]]
id = "run-tests-button"
title = "Run Tests"
icon = "play-circle"
command = "extension.runTests"
when = "hasTestFiles"

# Register a custom status bar item
[[ui_components.statusbar]]
id = "coverage-indicator"
text = "Coverage: ${coverage}%"
alignment = "right"
priority = 100
\end{lstlisting}

\subsection*{Dynamic Layout Modifications}
\addcontentsline{toc}{subsection}{Dynamic Layout Modifications}

Symphony's layout system is built on a flexible grid that extensions can modify dynamically:

\begin{infobox}[title=Adaptive Layout Philosophy]
Unlike traditional IDEs with fixed layouts, Symphony's UI adapts to the extensions installed and the user's workflow. An extension can request layout changes, and Symphony intelligently accommodates them while maintaining usability.
\end{infobox}

\subsubsection*{Layout Regions and Panels}

The UI framework defines several extensible regions:

\begin{compactlist}
    \item \textbf{Editor Area}: Central editing space with tab management
    \item \textbf{Sidebar Panels}: Left and right collapsible panels
    \item \textbf{Bottom Panel}: Terminal, output, and tool panels
    \item \textbf{Top Toolbar}: Action buttons and menu integration
    \item \textbf{Status Bar}: Information and quick actions
    \item \textbf{Modal Overlays}: Dialogs, command palette, notifications
\end{compactlist}

\subsubsection*{Dynamic Panel Management}

Extensions can create and manage panels dynamically:

\begin{lstlisting}[language=javascript, caption=Dynamic Panel Creation]
// Create a new panel programmatically
const panel = await symphony.ui.createPanel({
    id: 'debug-variables',
    title: 'Variables',
    location: 'sidebar-right',
    component: VariablesTreeComponent,
    data: debugSession.variables
});

// Show/hide panels based on context
symphony.workspace.onDidChangeActiveEditor((editor) => {
    if (editor.language === 'rust') {
        panel.show();
    } else {
        panel.hide();
    }
});

// Update panel content dynamically
debugSession.onVariablesChanged((variables) => {
    panel.updateData(variables);
});
\end{lstlisting}

\subsection*{Sidebars, Panels, and Overlays}
\addcontentsline{toc}{subsection}{Sidebars, Panels, and Overlays}

Symphony provides three primary UI extension mechanisms:

\subsubsection*{Sidebar Integration}

Sidebars provide persistent, contextual information and controls:

\begin{compactlist}
    \item \textbf{Tree Views}: File explorers, symbol outlines, test hierarchies
    \item \textbf{List Views}: Search results, problem lists, task queues
    \item \textbf{Form Views}: Settings panels, configuration interfaces
    \item \textbf{Custom Views}: Specialized interfaces for domain-specific tools
\end{compactlist}

\subsubsection*{Panel System}

Panels offer flexible workspace organization:

\begin{compactlist}
    \item \textbf{Tabbed Panels}: Multiple related views in a single container
    \item \textbf{Split Panels}: Side-by-side or stacked panel arrangements
    \item \textbf{Floating Panels}: Detachable panels for multi-monitor setups
    \item \textbf{Collapsible Panels}: Space-efficient panels that can be minimized
\end{compactlist}

\subsubsection*{Overlay System}

Overlays provide temporary, focused interactions:

\begin{compactlist}
    \item \textbf{Command Palette}: Quick action discovery and execution
    \item \textbf{Quick Open}: File and symbol navigation
    \item \textbf{Notifications}: Non-intrusive status updates
    \item \textbf{Modal Dialogs}: Configuration and confirmation interfaces
\end{compactlist}

\subsection*{Virtual DOM Bridge}
\addcontentsline{toc}{subsection}{Virtual DOM Bridge}

For extensions requiring maximum flexibility, Symphony provides a Virtual DOM bridge that allows direct manipulation of the UI:

\begin{successbox}
The Virtual DOM bridge enables extensions to create sophisticated, high-performance user interfaces while maintaining integration with Symphony's theming and accessibility systems.
\end{successbox}

\subsubsection*{Bridge Architecture}

The Virtual DOM bridge operates through several layers:

\begin{expandedlist}
    \item \textbf{React Integration}: Extensions can use React components with full lifecycle support
    \item \textbf{State Management}: Integration with Symphony's global state management
    \item \textbf{Event System}: Unified event handling across core and extension UI
    \item \textbf{Theme Integration}: Automatic application of user themes and preferences
    \item \textbf{Accessibility}: Built-in accessibility features and ARIA support
\end{expandedlist}

\subsubsection*{Performance Optimization}

The bridge includes several performance optimizations:

\begin{compactlist}
    \item \textbf{Virtual Scrolling}: Efficient rendering of large lists and trees
    \item \textbf{Lazy Loading}: Components load only when needed
    \item \textbf{Memoization}: Automatic caching of expensive computations
    \item \textbf{Batch Updates}: Efficient DOM updates through batching
\end{compactlist}

\subsection*{Component Registration and Lifecycle}
\addcontentsline{toc}{subsection}{Component Registration and Lifecycle}

Symphony manages the complete lifecycle of extension UI components:

\subsubsection*{Registration Process}

\begin{expandedlist}
    \item \textbf{Declaration}: Components declared in extension manifest
    \item \textbf{Validation}: Symphony validates component specifications
    \item \textbf{Registration}: Components registered with the UI system
    \item \textbf{Instantiation}: Components created when needed
    \item \textbf{Cleanup}: Automatic cleanup when extensions are disabled
\end{expandedlist}

\subsubsection*{Lifecycle Hooks}

Extensions can hook into various UI lifecycle events:

\begin{lstlisting}[language=javascript, caption=UI Component Lifecycle]
class CustomPanelComponent extends React.Component {
    componentDidMount() {
        // Register for workspace events
        this.subscription = symphony.workspace.onDidChangeFiles(
            this.handleFileChanges
        );
    }
    
    componentWillUnmount() {
        // Clean up subscriptions
        this.subscription.dispose();
    }
    
    handleFileChanges = (changes) => {
        // Update UI based on file system changes
        this.setState({ files: changes });
    }
}
\end{lstlisting}

\begin{alertbox}
The UI extensibility framework ensures that extensions can create rich, interactive interfaces while maintaining Symphony's performance characteristics and user experience consistency.
\end{alertbox}